\bookchapter{Глава 61. Быстрый старт со Swift — первый язык разработки для Apple}

\chapterintro{В этой главе мы всерьёз берёмся за язык. Не волнуйтесь — Swift славится «дружелюбием к новичкам», и зубрить учебник по синтаксису не нужно: освоив 20\% синтаксиса, вы сможете понимать и писать 80\% повседневного кода. Остальное будете подсматривать по ходу проекта. Запомните наш принцип обучения: понять важнее, чем запомнить; уметь пользоваться важнее, чем владеть в совершенстве.}

\booksection{Переменные и типы: даём данным имена}

Программа — это «данные + логика». Переменная — это имя, под которым данные хранятся. В Swift их объявляют через let и var:

\begin{lstlisting}[language=Swift]
let name = "Иван"        // let: константа (по умолчанию берём её)
var age = 18             // var: переменная (можно менять потом)

age = 19                 // <i class="fas fa-check-circle mr-1"></i> var можно изменить
// name = "Мария"         // <i class="fas fa-times-circle mr-1"></i> let изменить нельзя — ошибка компиляции

let score: Int = 100     // тип можно указать и явно
let pi = 3.14            // можно не писать — выведется сам

print("Привет, \(name)")   // строковая интерполяция: обратный слэш + скобки
\end{lstlisting}

Часто используемые типы — запомните одной таблицей:

\begin{center}
\begin{tabularx}{\linewidth}{|X|X|X|}
\hline
\textbf{Тип} & \textbf{Что хранит} & \textbf{Пример} \tabularnewline
\hline
Int & целые числа & 42 \tabularnewline
Double & дробные числа & 3.14 \tabularnewline
String & текст & "Привет" \tabularnewline
Bool & истина/ложь & true / false \tabularnewline
{}[Int] (массив) & набор данных & {}[1, 2, 3] \tabularnewline
\hline
\end{tabularx}
\end{center}

\begin{center}\small\color{gray}Часто используемые базовые типы Swift\end{center}

\begin{tcbinfo}{Совет: принцип «сначала let»}
Железное правило сообщества Swift: по умолчанию используйте let, а var — только когда это действительно нужно. Если там, где хватило бы let, вы написали var, Xcode предупредит вас жёлтым. Эта привычка избавит вас от целого класса багов «переменную случайно изменили» — чем меньше в коде меняется, тем легче о нём рассуждать.

\end{tcbinfo}

\booksection{Опциональные значения (Optional): «оберег от падений» в Swift}

Это важнейшее понятие Swift, и без всякой конкуренции. Одной фразой: если значение может «отсутствовать», Swift заставляет вас это обработать — язык на корню искореняет профессиональный кошмар «падение из-за нулевого указателя».

\begin{lstlisting}[language=Swift]
var name: String = "Иван"   // обычный тип: nil не допускается
// name = nil                // <i class="fas fa-times-circle mr-1"></i> ошибка компиляции — вас сразу остановят

var nickname: String? = nil  // добавляем вопросик: nil разрешён (опционал)
nickname = "Ваня"

// Опционал нужно «распаковать»:
if let n = nickname {        // if let: безопасная распаковка
    print("Привет, \(n)")
} else {
    print("Никнейма ещё нет")
}

let len = nickname?.count ?? 0   // ?? подставляет значение по умолчанию (оператор nil-coalescing)
print(len)
\end{lstlisting}

\begin{tcbwarn}{Внимание: «!» — опасная операция}
Написать nickname! (принудительное разворачивание) — всё равно что крикнуть компилятору «клянусь, там не nil»: ошибётесь — приложение рухнет на месте. Правило новичка: по возможности используйте if let / ?? и не используйте «!». Когда Xcode автодополняет код, он иногда подсовывает «!» — не забудьте это убрать.

\end{tcbwarn}

\booksection{Функции и замыкания: упаковываем логику в «кнопки»}

Функция — это поименованный кусок логики, как кнопка на пульте. Замыкание (Closure) — «функция без имени», которую можно передать другому одним движением руки:

\begin{lstlisting}[language=Swift]
// Функции: имя впереди, параметры и возвращаемый тип после
func greet(name: String) -> String {
    return "Привет, \(name)!"
}
print(greet(name: "Иван"))     // при вызове нужно указывать имя параметра

// Замыкания: определяем на месте, коротко и прямо
let numbers = [1, 2, 3, 4]
let doubled = numbers.map { $0 * 2 }   // $0 = текущий элемент
// [2, 4, 6, 8]

let big = numbers.filter { $0 > 2 }    // фильтрация
// [3, 4]
\end{lstlisting}

В SwiftUI замыкания повсюду: фигурные скобки после Button(“Нажми меня”) \{ ... \} — это и есть замыкание. Крепко запомните «комбо» для коллекций: map, filter, sorted — 90\% будущей обработки списков будет их комбинацией.

\booksection{Структуры и протоколы: «лего»-система Swift}

Официальное предпочтение Swift ясно: сначала структуры (struct), а классы (class) — только когда нужно наследование. Структура — «тип значения»: при копировании каждая живёт своей жизнью, и это само по себе избавляет от целого класса багов:

\begin{lstlisting}[language=Swift]
// Структура: коробка для данных
struct Record {
    var id = UUID()
    var amount: Double
    var category: String
    var date: Date

    // У структуры есть встроенный «инициализатор по членам»: создавать можно сразу
}
let r = Record(amount: 25.5, category: "Еда", date: Date())

// Протокол: договорённость «какие способности обязательны»
protocol Identifiable {
    var id: UUID { get }
}
struct Record: Identifiable { ... }   // соответствие протоколу = обещание выполнить его требования
\end{lstlisting}

В SwiftUI почти всё — struct (View — это протокол). Если сейчас фраза непонятна — не страшно, в главе 62 вы будете встречаться с ней каждый день. В мире Swift сперва обнимите struct.

\booksection{Шпаргалка по синтаксису для начинающих с нуля}

\begin{center}
\begin{tabularx}{\linewidth}{|X|X|}
\hline
\textbf{Что нужно сделать} & \textbf{Как пишется} \tabularnewline
\hline
Объявить константу/переменную & \texttt{let x = 1}/\texttt{var x = 1} \tabularnewline
Интерполяция строк & "Привет, \textbackslash{}(name)" \tabularnewline
Объявить функцию & func имя(параметр: Тип) -> Тип Результата \{ ... \} \tabularnewline
Условие & if условие \{ ... \} else \{ ... \} \tabularnewline
Перебор & for x in массив \{ ... \} \tabularnewline
Безопасное разворачивание & if let x = опциональное Значение \{ ... \} \tabularnewline
Значение по умолчанию & x ?? значение По Умолчанию \tabularnewline
Объявить структуру & \texttt{struct X \{ var a: Int; var b: String \}} \tabularnewline
Замыкание & \texttt{\{ \$0 * 2 \}} \tabularnewline
\hline
\end{tabularx}
\end{center}

\begin{center}\small\color{gray}Шпаргалка по синтаксису Swift (для новичков)\end{center}

\begin{tcbpractice}{Попробуйте сами: 15 минут игры со Swift в Playgrounds}
Есть Mac — откройте Xcode Playground; нет Mac — используйте Swift Playgrounds на iPad или онлайн-песочницу. Прогоните каждый кусочек кода из этой главы, меняйте числа и текст, нарочно наделайте ошибок и посмотрите, как на них реагирует Xcode. Эти 15 минут «сначала поиграть, потом учиться» сделают старт главы 62 вдвое легче.

\end{tcbpractice}

\begin{tcbsummary}{Краткие итоги главы}
\begin{itemize}
  \item let — по умолчанию, var — с осторожностью; типы выводятся автоматически, но можно указать и явно после двоеточия
  \item Опциональные значения — душа Swift: если значение может быть пустым, это обязательно обрабатывается; if let / ?? — оберег, «!» — опасная операция
  \item При вызове функции указываются имена параметров; замыкание = функция без имени, в SwiftUI оно повсюду
  \item Приоритет — struct, соблюдайте протоколы (Protocol); в SwiftUI почти всё — struct
  \item 20\% синтаксиса делают 80\% работы: сначала научитесь пользоваться, потом совершенствуйтесь
\end{itemize}

\end{tcbsummary}
